

\section{The Art and Science of Data Visualization through Plotting }

In this session, we shift fully into an \textbf{Iris-only} workflow. There are \textbf{no robots} in this tutorial.
Recall from Session 8 that we already produced NumPy arrays (\texttt{X}, \texttt{y}) and class centroids.
Now, building on Session 8's outputs, we visualize classifier behavior and experiment results.

\noindent
\textbf{Deliverable:} \texttt{session9/session9\_iris\_template.py}

\noindent
You will work on two key tasks:
\begin{enumerate}\setlength{\itemsep}{2pt}
\item \textbf{Systematic parameter experiments:} run repeated experiments by changing a classifier setting (e.g., distance metric, threshold, or feature subset), and record the accuracy.
\item \textbf{Visualization:} plot how performance changes across settings, using box plots and line plots.
\end{enumerate}

\noindent
This tutorial is important because good data practice means:
\begin{itemize}\setlength{\itemsep}{0pt}
\item running \textbf{multiple trials} to reduce randomness and draw reliable conclusions,
\item saving results in a structured form (\textbf{list of dictionaries} $\rightarrow$ \textbf{DataFrame}),
\item visualizing distributions (box plot) and trends (line plot).
\end{itemize}

\subsection{Prerequisites}
You need:
\begin{itemize}\setlength{\itemsep}{0pt}
\item \texttt{numpy}
\item \texttt{pandas}
\item \texttt{matplotlib}
\item an \texttt{iris.csv} file (same dataset used throughout the module)
\end{itemize}

\noindent
Install if needed:
\begin{verbatim}
pip install numpy pandas matplotlib openpyxl
\end{verbatim}

\newpage
\subsection{Calculating Accuracy for Iris Experiments}

In session 7, you built an Iris classifier using inheritance (e.g., a base classifier and child classifiers for different species).
In session 8, you computed \textbf{centroids} using NumPy.

\noindent
In session 9, we combine these ideas into an \textbf{experiment pipeline}:
\begin{itemize}\setlength{\itemsep}{0pt}
\item load Iris into NumPy arrays,
\item run a classifier,
\item compute accuracy,
\item repeat under different settings,
\item store and visualize results.
\end{itemize}

\subsubsection{Your Task: Reuse session 8 Modules (Iris-only)}

Reuse your session 8 modules directly instead of rewriting them.
Your folder structure should look like this:

\begin{verbatim}
session8/
|__ iris_session8_utils.py
|__ iris_session8_models.py
|__ session8_iris_template.py

session9/
|__ session9_iris_template.py
|__ session9_plot_accuracy.py
\end{verbatim}

\noindent
\textbf{Important:} Session 9 should import from Session 8 (for example \texttt{session8.iris\_session8\_utils} and \texttt{session8.iris\_session8\_models}) so continuity is preserved.

\newpage
\subsection{Refactoring for Flexibility}
\label{sec:activity2}

In this activity, we refactor our experiment runner so that key settings are not hardcoded.
We will focus on three parameters:
\begin{itemize}\setlength{\itemsep}{0pt}
\item \texttt{test\_size} (fraction of data used for testing),
\item \texttt{seed} (controls randomness),
\item \texttt{metric} (distance metric: \texttt{"l2"} or \texttt{"l1"}).
\end{itemize}

\subsubsection{Identifying the Need for Refactoring}

A common beginner mistake is hardcoding values:
\begin{minted}[fontsize=\small, linenos, breaklines=true]{python}
# BAD PRACTICE (example)
test_size = 0.2
seed = 1
metric = "l2"
\end{minted}

\noindent
This is problematic because:
\begin{itemize}\setlength{\itemsep}{0pt}
\item you must edit code repeatedly to test new settings,
\item it is hard to run systematic experiments,
\item results become difficult to reproduce.
\end{itemize}

\subsubsection{Your Task: Parameterize the Experiment Runner and Add Return Value}
\label{subsubsec:activity_parametrize_main}

Create (or edit) \texttt{session9/session9\_iris\_template.py}. We will define a function called \texttt{run\_experiment} that:
\begin{itemize}\setlength{\itemsep}{0pt}
\item accepts parameters,
\item runs one train-test split,
\item returns \textbf{accuracy}.
\end{itemize}

\noindent
\textbf{Target function signature:}
\begin{minted}[fontsize=\small, linenos, breaklines=true]{python}
def run_experiment(test_size=0.2, seed=1, metric="l2"):
    ...
    return accuracy
\end{minted}

\newpage
\subsection{A Simple Iris Centroid Classifier (Built-in Option)}
\label{sec:simple_centroid_classifier}

If your session 7 inheritance-based classifier is not ready, use this centroid-based classifier.
It stays Iris-only and uses only NumPy.

\subsubsection{Reuse \texttt{session8/iris\_session8\_models.py} (recommended)}
\noindent
If your session 8 model interface is ready, import it directly in Session 9:
\begin{minted}[fontsize=\small, linenos, breaklines=true]{python}
from session8.iris_session8_models import NumpyCentroidClassifier
\end{minted}

\subsubsection{Fallback: Create \texttt{iris\_model.py}}
\begin{minted}[fontsize=\small, linenos, breaklines=true]{python}
# iris_model.py
import numpy as np

class CentroidClassifier:
    """
    A simple centroid classifier:
    - compute class centroids from training data
    - predict by nearest centroid
    metric: "l2" (euclidean) or "l1" (manhattan)
    """

    def __init__(self, metric="l2"):
        self.metric = metric
        self.centroids = {}
        self.classes_ = None

    def fit(self, X_train, y_train):
        self.classes_ = np.unique(y_train)
        for c in self.classes_:
            self.centroids[c] = np.mean(X_train[y_train == c], axis=0)

    def _dist(self, a, b):
        if self.metric == "l2":
            return np.sqrt(np.sum((a - b) ** 2))
        elif self.metric == "l1":
            return np.sum(np.abs(a - b))
        else:
            raise ValueError("metric must be 'l2' or 'l1'")

    def predict_one(self, x):
        best_class = None
        best_dist = float("inf")
        for c in self.classes_:
            d = self._dist(x, self.centroids[c])
            if d < best_dist:
                best_dist = d
                best_class = c
        return best_class

    def predict(self, X):
        return np.array([self.predict_one(x) for x in X], dtype=str)
\end{minted}

\newpage
\subsection{Activity 3: Implement Train/Test Split and Accuracy (NumPy-only)}

\subsubsection{Create \texttt{session9/iris\_eval.py}}
\begin{minted}[fontsize=\small, linenos, breaklines=true]{python}
# iris_eval.py
import numpy as np

def train_test_split(X, y, test_size=0.2, seed=1):
    """
    NumPy-only train/test split.
    """
    n = X.shape[0]
    rng = np.random.default_rng(seed)
    idx = np.arange(n)
    rng.shuffle(idx)

    test_n = int(n * test_size)
    test_idx = idx[:test_n]
    train_idx = idx[test_n:]

    return X[train_idx], X[test_idx], y[train_idx], y[test_idx]

def accuracy_score(y_true, y_pred):
    return float(np.mean(y_true == y_pred))
\end{minted}

\newpage
\subsection{Activity 4: Run One Experiment and Print an Informative Output}
\label{sec:activity4_one_run}

\subsubsection{Update \texttt{session9/session9\_iris\_template.py}}
\begin{minted}[fontsize=\small, linenos, breaklines=true]{python}
# session9/session9_iris_template.py
import numpy as np
from session8.iris_session8_utils import load_iris_csv
from session9.iris_eval import train_test_split, accuracy_score

# Option A: your session 7 model (if you have it)
# from iris_classifier_session7 import YourClassifier

# Option B: centroid classifier (provided)
from session8.iris_session8_models import NumpyCentroidClassifier

def run_experiment(test_size=0.2, seed=1, metric="l2"):
    X, y = load_iris_csv("iris.csv", skip_header=1)
    X_train, X_test, y_train, y_test = train_test_split(X, y, test_size=test_size, seed=seed)

    model = NumpyCentroidClassifier(metric=metric)
    model.fit(X_train, y_train)
    y_pred = model.predict(X_test)

    acc = accuracy_score(y_test, y_pred)
    return acc

if __name__ == "__main__":
    acc = run_experiment(test_size=0.2, seed=1, metric="l2")
    print(f"One run | test_size=0.2 | seed=1 | metric=l2 | accuracy={acc:.3f}")
\end{minted}

\subsubsection{Your Task}
\begin{enumerate}\setlength{\itemsep}{2pt}
\item Run \texttt{session9/session9\_iris\_template.py}.
\item Change \texttt{metric} to \texttt{"l1"} and run again.
\item Observe the accuracy difference.
\end{enumerate}

\newpage
\subsection{Running Multiple Experiments Dynamically}

Now we want to run multiple experiments with different parameter settings.
For example:
\begin{itemize}\setlength{\itemsep}{0pt}
\item test \texttt{metric} in \{\texttt{"l1"}, \texttt{"l2"}\}
\item test multiple \texttt{seed} values
\item test multiple \texttt{test\_size} values
\end{itemize}

\subsubsection{Your Task: Implement Informative and Looped Output}
\label{subsubsec:activity_informative_loop_output}

Add this function inside \texttt{session9/session9\_iris\_template.py}:

\begin{minted}[fontsize=\small, linenos, breaklines=true]{python}
def run_single_setting_sweep(metrics, seeds, test_size=0.2):
    print(f"--- Sweep: test_size={test_size} ---")
    for metric in metrics:
        for seed in seeds:
            acc = run_experiment(test_size=test_size, seed=seed, metric=metric)
            print(f"metric={metric} | seed={seed} | acc={acc:.3f}")
\end{minted}

\noindent
Then call it in the main block:
\begin{minted}[fontsize=\small, linenos, breaklines=true]{python}
if __name__ == "__main__":
    run_single_setting_sweep(metrics=["l1", "l2"], seeds=[1, 2, 3, 4, 5], test_size=0.2)
\end{minted}

\newpage
\subsection{Good Scientific Practice by Running Multiple Trials and Informative Output}
\label{sec:activity3_multipletrials}

In this activity, we run \textbf{multiple trials} for each configuration, and store results.

\subsubsection{Your Task: Implement \texttt{run\_multiple\_trials}}
\label{subsubsec:implement_run_multiple_trials}

\begin{minted}[fontsize=\small, linenos, breaklines=true]{python}
def run_multiple_trials(metrics, test_sizes, num_trials):
    """
    Returns a list of dicts:
    {'metric':..., 'test_size':..., 'trial':..., 'seed':..., 'accuracy':...}
    """
    results_list = []

    for metric in metrics:
        for test_size in test_sizes:
            print(f"\n[Config] metric={metric} | test_size={test_size}")
            for trial in range(1, num_trials + 1):
                seed = trial  # simple choice: seed=1..num_trials
                acc = run_experiment(test_size=test_size, seed=seed, metric=metric)

                print(f"  Trial {trial}/{num_trials} | seed={seed} | acc={acc:.3f}")

                run_data = {
                    "metric": metric,
                    "test_size": test_size,
                    "trial": trial,
                    "seed": seed,
                    "accuracy": acc
                }
                results_list.append(run_data)

    return results_list
\end{minted}

\subsubsection{Main block call}
\begin{minted}[fontsize=\small, linenos, breaklines=true]{python}
if __name__ == "__main__":
    all_results_data = run_multiple_trials(
        metrics=["l1", "l2"],
        test_sizes=[0.2, 0.3, 0.4],
        num_trials=10
    )
    print("\n--- Raw Results (list of dicts) ---")
    print(all_results_data[:5])
\end{minted}

\newpage
\subsection{Analyzing Results with Pandas and Saving Data}
\label{sec:saving_as_pandas}

Now we convert the list of dictionaries into a DataFrame and save it.

\subsubsection*{Your Task: Use Pandas to Process and Save Results}
\label{subsubsec:activity5_task}

At the top of \texttt{session9/session9\_iris\_template.py}, add:
\begin{minted}[fontsize=\small, linenos]{python}
import pandas as pd
\end{minted}

\noindent
Then after \texttt{all\_results\_data} is created:
\begin{minted}[fontsize=\small, linenos, breaklines=true]{python}
results_df = pd.DataFrame(all_results_data)
print(results_df.head())

output_filename = "iris_experiment_results.xlsx"
results_df.to_excel(output_filename, index=False)
print(f"Saved: {output_filename}")
\end{minted}

\newpage
\subsection{Analyzing and Visualizing Results}
\label{sec:activity6_plotting}

Once your results are in a DataFrame, we can visualize performance patterns.

\subsubsection{Basic Analysis: Grouping and Aggregating}

\begin{minted}[fontsize=\small, linenos, breaklines=true]{python}
avg_acc = results_df.groupby(["metric", "test_size"])["accuracy"].mean().reset_index()
print(avg_acc)
\end{minted}

\subsubsection{Visualization with Matplotlib}

Create a new file: \texttt{session9/session9\_plot\_accuracy.py}

\begin{minted}[fontsize=\small, linenos, breaklines=true]{python}
import pandas as pd
import matplotlib.pyplot as plt

df = pd.read_excel("iris_experiment_results.xlsx")

# Box plot: accuracy distribution by metric
df.boxplot(column="accuracy", by="metric")
plt.title("Accuracy Distribution by Distance Metric")
plt.suptitle("")
plt.xlabel("Metric")
plt.ylabel("Accuracy")
plt.grid(True)
plt.show()

# Line plot: mean accuracy vs test_size (separate line per metric)
avg = df.groupby(["metric", "test_size"])["accuracy"].mean().reset_index()

for metric in sorted(avg["metric"].unique()):
    sub = avg[avg["metric"] == metric]
    plt.figure()
    plt.plot(sub["test_size"], sub["accuracy"], marker="o")
    plt.title(f"Mean Accuracy vs Test Size ({metric})")
    plt.xlabel("Test Size")
    plt.ylabel("Mean Accuracy")
    plt.grid(True)
    plt.show()
\end{minted}

\subsubsection{Your Task}
\begin{enumerate}\setlength{\itemsep}{2pt}
\item Run \texttt{session9/session9\_iris\_template.py} to generate \texttt{iris\_experiment\_results.xlsx}.
\item Run \texttt{session9/session9\_plot\_accuracy.py}.
\item Screenshot your plots for submission.
\end{enumerate}

\newpage
\subsection{Optional: LLM for Code Understanding (Iris-only)}

\noindent
Pick an LLM tool (ChatGPT / Claude / Gemini / etc.). Paste \texttt{iris\_model.py} and ask:

\begin{quote}
    Explain how NumPy is used in this classifier. For each NumPy line, tell me the input shape and output shape. Then explain why L1 vs L2 distance might change predictions.
\end{quote}

\noindent
Write 3--5 lines (as comments in \texttt{session9/session9\_iris\_template.py}):
\begin{itemize}\setlength{\itemsep}{0pt}
\item One thing the LLM clarified,
\item One thing you still find confusing.
\end{itemize}

\subsection{Submission}
Submit:
\begin{itemize}\setlength{\itemsep}{0pt}
\item \texttt{session9/session9\_iris\_template.py}
\item \texttt{session9/session9\_plot\_accuracy.py}
\item \texttt{session8/iris\_session8\_utils.py}
\item \texttt{session8/iris\_session8\_models.py}
\item \texttt{session9/iris\_eval.py}
\item \texttt{iris\_experiment\_results.xlsx}
\item screenshots of plots (box plot + line plot)
\end{itemize}

\begin{center}
{\color{red}\Large Reminder: Iris-only tutorial. No robot simulation content.}
\end{center}
